\documentclass[11pt]{article}
\usepackage[a4paper,margin=25mm]{geometry}
\usepackage{amsmath,amssymb,booktabs,array}
\usepackage[authoryear,round]{natbib}
\usepackage[hidelinks]{hyperref}
\usepackage{microtype}
\setlength{\emergencystretch}{1.5em}
\setlength{\parindent}{0pt}
\setlength{\parskip}{5pt}
\newcommand{\E}{\mathbb E}
\newcommand{\Var}{\operatorname{Var}}
\newcommand{\Cov}{\operatorname{Cov}}
\newcommand{\diag}{\operatorname{diag}}
\newcommand{\sgn}{\operatorname{sign}}
\title{Online Resource 1: Mathematical checks of\protect\\ reward-to-update claims}
\author{Yanqi Li\\University of Southern Queensland\\\texttt{Yanqi.Li@unisq.edu.au}}
\date{}
\begin{document}
\begin{center}
{\large\bfseries Online Resource 1: Mathematical checks of reward-to-update claims}\par
\vspace{7pt}
Yanqi Li\\University of Southern Queensland\\\texttt{Yanqi.Li@unisq.edu.au}
\end{center}

Supplement to \emph{Enhancing Reinforcement Learning with Large Language Models: A Survey of Grounded Reward, Policy, and Environment Interfaces}, submitted to \emph{Artificial Intelligence Review}.

This note gives the complete constructions behind the survey's two checks: preservation of discounted task preferences, and the passage from reward quality to a policy update. It connects each construction to a published formula in ARM-FM, SWE-RM, or ProCredit. The numerical results below are analytical examples constructed for the survey; they are not measurements from those papers' training runs. Two accompanying standard-library Python programs check the constructions: one reproduces the numerical examples by exact arithmetic and exhaustive enumeration; the other verifies the calibrated family using exact identities, fixed-seed numerical checks, finite differences, and independent group enumeration. Neither program trains a policy.

\section{Discounted preferences under ARM-FM}
\label{sec:arm}

\paragraph{Source condition and learning target.}
Appendix A.5 of \citet{Castanyer2026ARMFM} states optimal-policy preservation when a reward machine has no positive-reward cycles and its terminal reward exceeds the cumulative reward of every nonterminal trajectory. Its DQN implementation adds the task and machine rewards. With environment state $s$, machine state $u$, language embedding $\phi(u)$, next-state quantities $s',u'$, terminal indicator $d$, and target network $\widehat Q$, Algorithm~1, line~20 (PDF p.~19), gives
\begin{equation}
y=r_{\rm task}+r_{\rm RM}
  +\gamma(1-d)\max_{a'}\widehat Q(s',\phi(u'),a').
\label{eq:dqn}
\end{equation}

\paragraph{Construction.}
Consider a deterministic episodic task with initial state $s$, intermediate state $x$, and terminal goal $g$. One action reaches $g$ immediately; another reaches $x$, from which the next action reaches $g$. The task pays one on entry to $g$ and zero otherwise. Its discount is $\gamma=.99$.

The reward machine accurately recognizes entry to $x$ and task completion. It accepts completion on either route; the intermediate event is optional for this task. Such optional events are permitted by the published formalism: the Meta-World machine includes both a direct $u_0\to u_2$ grasp transition and a route through $u_1$ that earns an additional near-object reward \citep[Appendix A.9, PDF p.~32]{Castanyer2026ARMFM}. The constructed machine pays $b=.2$ at $x$ and one at $g$. There are no positive-reward cycles; unused events can have zero-reward self-loops. Every nonterminal reward total is at most $.2$, strictly below the terminal reward one. Both conditions in the source are satisfied.

\begin{center}
\renewcommand{\arraystretch}{1.15}
\begin{tabular}{lrrr}
\toprule
Path & Task return & Machine return & Additive return\\
\midrule
$s\to g$ & $1$ & $1$ & $2$\\
$s\to x\to g$ & $.99$ & $.2+.99=1.19$ & $.2+2(.99)=2.18$\\
\bottomrule
\end{tabular}
\end{center}

The task prefers immediate completion, whereas the machine-only and additive objectives prefer the delayed path. At the Bellman fixed point of Eq.~\eqref{eq:dqn}, the terminal target is two. The delayed action therefore has value $.2+.99(2)=2.18$, above the direct action's value two. This is a comparison of the target's fixed-point preferences, not the first network update from an arbitrary initialization.

For a general intermediate bonus $b$, the machine-only reversal requires $b>1-\gamma$; the additive reversal requires
\begin{equation}
b>2(1-\gamma).
\label{eq:rmthreshold}
\end{equation}
The source's requirement $b<1$ does not exclude either. Both constructed paths succeed. The counterexample concerns discounted preference preservation, not a failed trajectory outranking completion, and it does not establish failure in the source's experimental tasks.

\section{Calibration and a policy-gradient direction}
\label{sec:calibration}

\subsection{The missing conditional covariance}
For a trajectory $\tau$, let $r(\tau)$ be a fixed predicted score, $c(\tau)\in\{0,1\}$ its correctness, and
$h(\tau)=\nabla_\theta\log\pi_\theta(\tau)$ its policy score function. Define $J_f(\theta)=\E_{\pi_\theta}[f]$. Appendix C.3 of \citet{Shum2026SWE_RM} states calibration without a uniform-in-policy qualification. At a fixed sampling policy, this condition is
\begin{equation}
\E[c\mid r]=r
\label{eq:calibration}
\end{equation}
From this condition, the source infers equality of $\E[rh]$ and $\E[ch]$ (source Eq.~12, PDF p.~16). Under the usual differentiability conditions for a fixed evaluator,
\begin{align}
\nabla J_r-\nabla J_c
 &=\E[(r-c)h] \nonumber\\
 &=-\E[\Cov(c,h\mid r)].
\label{eq:covariance}
\end{align}
To see the second equality, condition on $r$ and use
$\E[ch\mid r]=\E[c\mid r]\E[h\mid r]+\Cov(c,h\mid r)$.
Calibration fixes the average correctness within a score group. It does not force the policy sensitivities of correct and incorrect trajectories in that group to agree.

\subsection{An exactly calibrated three-action example}
Consider a one-step task with three actions and independent softmax logits:
\[
\pi_z(i)=\frac{\exp(z_i)}{\sum_j\exp(z_j)},\qquad z_i^0=\log p_i.
\]
Use the following probabilities, correctness labels, and scores:
\begin{center}
\begin{tabular}{cccc}
\toprule
Action $i$ & $p_i$ & $c_i$ & $r_i$\\
\midrule
1 & $.1$ & $1$ & $1$\\
2 & $.6$ & $1$ & $2/3$\\
3 & $.3$ & $0$ & $2/3$\\
\bottomrule
\end{tabular}
\end{center}
The evaluator is exactly calibrated: the score-one action is correct, and correctness in the lower-score group is $.6/(.6+.3)=2/3$. Both expected score and success equal $.7$.

For fixed values $f_i$,
\begin{equation}
\left.\partial_{z_i}J_f\right|_{z^0}=p_i(f_i-J_f).
\label{eq:logit}
\end{equation}
Thus the success and score gradients are
\begin{equation}
g_c=(.03,.18,-.21),\qquad g_r=(.03,-.02,-.01).
\label{eq:gcgr}
\end{equation}
Along an ordinary Euclidean logit step $z(\eta)=z^0+\eta g_r$,
\begin{equation}
\left.\frac{dJ_r}{d\eta}\right|_0=\lVert g_r\rVert^2=.0014,
\qquad
\left.\frac{dJ_c}{d\eta}\right|_0=g_c^\top g_r=-.0006.
\label{eq:local}
\end{equation}
A sufficiently small positive step increases score and decreases success. No later distribution shift is needed to obtain this local discrepancy. Calibration throughout a neighborhood of policies would instead imply equality of the two objectives there, and hence equality of their gradients. That is stronger than Eq.~\eqref{eq:calibration} at one policy.

\subsection{A family of calibrated base policies}
\label{sec:calibration_family}
For $a,b,d>0$ with $a+b+d=1$, write $m=1-a=b+d$ and set
\begin{equation}
 p=(a,b,d),\qquad c=(1,1,0),\qquad
 r=\left(1,\frac{b}{m},\frac{b}{m}\right).
 \label{eq:family_setup}
\end{equation}
The score-one action is correct, and correctness conditional on the lower
score is $b/m$, so the evaluator is exactly calibrated at $p$.
For each chosen base point, freeze these scores and define
$J_f(z)=\sum_i\pi_z(i)f_i$, where
$\pi_z(i)=e^{z_i}/\sum_j e^{z_j}$ and $z_i^0=\log p_i$.
The gradients below are evaluated in the full three-logit space at $z^0$.
They do not differentiate the rule that constructs $r$ from the base point.

Since $J_c(z^0)=J_r(z^0)=1-d$ and
$\partial_{z_i}J_f=p_i(f_i-J_f)$, direct calculation gives
\begin{align}
 g_c&=d\bigl(a,b,-(1-d)\bigr),\nonumber\\
 g_r&=ad\left(1,-\frac b m,-\frac d m\right),\label{eq:family_gradients}\\
 g_c-g_r&=\frac{bd}{m}(0,1,-1).\nonumber
\end{align}
Thus the gradient equality fails at every interior point in this family,
including points where the two gradients still have a positive inner product.
The stronger condition of a negative success derivative is
\begin{align}
 g_c^\top g_r
 &=\frac{ad^2}{m}\left[am-b^2+d(1-d)\right]\nonumber\\
 &=\frac{ad^2}{1-a}B(a,d),\label{eq:family_dot}\\
 B(a,d)&=(1-a)(2a-1)+(3-2a)d-2d^2.\nonumber
\end{align}
The prefactor is positive in the interior. Also $g_r\ne0$, so an ordinary
Euclidean step $z^0+\eta g_r$ has
$\dot J_r(0)=\lVert g_r\rVert^2>0$ and
$\dot J_c(0)<0$ precisely when $B(a,d)<0$.
Smoothness then gives the corresponding strict changes for all sufficiently
small positive step sizes.

\paragraph{The complete strict negative region.}
For fixed $a$, $B$ is a concave quadratic in $d$ on $[0,1-a]$.
Its endpoint values are $(1-a)(2a-1)$ and $2a(1-a)$.
For $a\geq1/2$, both endpoints are nonnegative and the interior is positive.
For $0<a<1/2$, the lower root lies inside this interval and the upper root
lies beyond it. Hence the complete interior region is
\begin{equation}
 0<a<\frac12,\qquad
 0<d<d_*(a):=
 \frac{3-2a-\sqrt{1+12a(1-a)}}4.
 \label{eq:family_negative_region}
\end{equation}
At $a=1/10$,
\begin{equation}
 d_*(1/10)=\frac{7-\sqrt{13}}{10}
 =0.3394448724536011\ldots.
 \label{eq:family_slice_root}
\end{equation}
For example, the open rectangle
$a\in(.09,.11)$, $d\in(.29,.31)$ lies wholly in the negative region.
Both partial derivatives of $B$ are positive there, while its supremum
at the upper corner is $B(.11,.31)=-.0246$.
This establishes a set of positive two-dimensional measure in the
base-policy construction, rather than an isolated numerical example.
Across this two-dimensional set, the selected score $b/(1-a)$ can change.
It is not a claim that one fixed evaluator remains calibrated on an open
neighborhood of the full policy simplex, nor an estimate of how often a
trained system encounters this region.

\paragraph{A fixed evaluator and a simpler parametrization.}
Put $q=b/(1-a)$, so that
$p=(a,(1-a)q,(1-a)(1-q))$ and $r=(1,q,q)$.
Equation~\eqref{eq:family_dot} becomes
\begin{equation}
 g_c^\top g_r
 =a(1-a)^2(1-q)^2
 \left[2a(q^2-q+1)-q(2q-1)\right].
 \label{eq:family_q_dot}
\end{equation}
Consequently, its strict negative region can equivalently be written as
\begin{equation}
 \frac12<q<1,\qquad
 0<a<\frac{q(2q-1)}{2(q^2-q+1)}.
 \label{eq:family_q_region}
\end{equation}
In particular, fixing $q=2/3$ fixes the evaluator throughout the family and
reduces the condition to $0<a<1/7$:
\begin{equation}
 g_c^\top g_r=\frac{2a(1-a)^2(7a-1)}{81},\qquad
 \lVert g_r\rVert^2=\frac{14a^2(1-a)^2}{81}.
 \label{eq:family_fixed_q}
\end{equation}
The numerical example above corresponds to $a=.1$.

The distinction between the base-policy family and the update coordinates is
essential. Along the curve with $q$ fixed,
$J_c(p(a))=J_r(p(a))=q+a(1-q)$, so both derivatives with respect to $a$
are $1-q$. Equivalently,
\[
 \frac{dz^0}{da}=\left(\frac1a,-\frac1{1-a},-\frac1{1-a}\right),\qquad
 g_c^\top\frac{dz^0}{da}=g_r^\top\frac{dz^0}{da}=1-q.
\]
The unfavorable direction in Eq.~\eqref{eq:family_fixed_q} is the ambient
ordinary logit gradient, which need not be tangent to this calibrated curve.
For example, it changes the low-score odds at rate
$\dot{\log(p_2/p_3)}=(g_r)_2-(g_r)_3=-ad(2q-1)$.
For the negative region, $q>1/2$, so those odds move against the correct
low-score action even though the two low-score actions have identical
predicted scores.

\section{Connecting the family to SWE-RM's GSPO surrogate}
\label{sec:gspo}

\subsection{The initial derivative}
SWE-RM uses group sequence policy optimization (GSPO), a sequence-normalized ratio, and standardized group advantages \citep[Appendix E.1, Eqs.~14--15, PDF pp.~20--21]{Shum2026SWE_RM}. For a group of $G$ trajectories from the behavior policy, write
\begin{equation}
\rho_i(\theta)=
\left[\frac{\pi_\theta(\tau_i)}{\pi_{\rm old}(\tau_i)}\right]^{1/L_i},
\qquad
A_i=\frac{r_i-\bar r}{s_G},
\label{eq:gsporatio}
\end{equation}
where $L_i$ is the positive trajectory-length factor in the source ratio,
$\bar r=G^{-1}\sum_i r_i$, and $s_G$ is the group standard deviation. The clipped surrogate averages
$\min\{\rho_i A_i,\operatorname{clip}(\rho_i,1-\varepsilon,1+\varepsilon)A_i\}$ over the group. At $\theta=\theta_{\rm old}$, the ratio is one and clipping is inactive for a clipping interval of positive width. Holding the sampled advantages and behavior policy fixed gives
\begin{equation}
\nabla_\theta\widehat J
=\frac1G\sum_{i=1}^G\frac{A_i}{L_i}
\nabla_\theta\log\pi_\theta(\tau_i).
\label{eq:gspoinitial}
\end{equation}
The source reports no KL loss and uses 16 rollouts per problem. The calculation below concerns this initial surrogate derivative with no additional loss gradient; it does not assume a particular optimizer implementation.

\subsection{Equal-length groups over the calibrated family}
\label{sec:family_gspo}
Use the family in Section~\ref{sec:calibration_family}, with $q=b/(1-a)$, one generated token per trajectory, and $G\geq2$ members sampled independently from the base policy. For this calculation let
$A_i=(r_i-\bar r)/(s_G+\epsilon_s)$ with $\epsilon_s>0$.
The source displays $s_G$ without specifying a stabilizer or the treatment of constant-score groups. The positive stabilizer used here makes those groups contribute zero without an undefined division.

Let $N\sim\operatorname{Binomial}(G,a)$ count score-one samples, $u=N/G$, and $\Delta=1-q>0$. Take $t_G=1$ for population standard deviation or $t_G=\sqrt{G/(G-1)}$ for sample standard deviation. The high- and low-score centered rewards are $\Delta(1-u)$ and $-\Delta u$, with common denominator
$t_G\Delta\sqrt{u(1-u)}+\epsilon_s$.
Since group advantages sum to zero, the common $-p$ terms in the softmax scores $e_i-p$ cancel. Conditional on $N$, the low-score actions occur in expected proportions $q$ and $1-q$. Therefore
\begin{align}
\E[\widehat g\mid N]&=k(u)(1,-q,-(1-q)),\nonumber\\
k(u)&=\frac{\Delta u(1-u)}{t_G\Delta\sqrt{u(1-u)}+\epsilon_s},\nonumber\\
\E[\widehat g]&=K_{G,\epsilon_s}(a,q)(1,-q,-(1-q)),\label{eq:family_gspo_direction}\\
K_{G,\epsilon_s}(a,q)&=\sum_{n=0}^G {G\choose n}a^n(1-a)^{G-n}k(n/G)>0.\nonumber
\end{align}
Here $\widehat g=\nabla_z\widehat J$. Strict positivity follows because $0<a<1$ and $G\geq2$ give positive probability to a mixed-score group. From Eq.~\eqref{eq:family_gradients},
\begin{equation}
\E[\widehat g]=\frac{K_{G,\epsilon_s}(a,q)}{ad}g_r.
\label{eq:family_gspo_collinear}
\end{equation}
Standardization thus changes the magnitude, but not the initial expected logit direction, throughout this two-score family. In particular, for the fixed evaluator $q=2/3$,
\begin{equation}
g_c^\top\E[\widehat g]
=\frac{2(1-a)(7a-1)}{27}K_{G,\epsilon_s}(a,2/3)<0
\quad\text{if }0<a<1/7.
\label{eq:family_gspo_fixed_q}
\end{equation}

Removing the standard-deviation divisor leaves $A_i=r_i-\bar r$. The self-included group mean gives
\begin{equation}
\E[\widehat g_{\rm centered}]=\frac{G-1}{G}g_r,
\qquad
\dot J_c=\frac{G-1}{G}g_c^\top g_r.
\label{eq:sweno}
\end{equation}
Indeed, write $\bar r=(r_i+\sum_{j\ne i}r_j)/G$. Independence makes the other members' contribution to the expected score-function product zero, leaving the coefficient $(G-1)/G$. The entire strict negative region in Eq.~\eqref{eq:family_q_region} therefore persists after removing standardization.

\paragraph{The original numerical example as a limit.}
For $a=.1$, $q=2/3$, and population standard deviation, the limit $\epsilon_s\downarrow0$ gives
\begin{equation}
\kappa_G:=\lim_{\epsilon_s\downarrow0}K_{G,\epsilon_s}(.1,2/3)
=\sum_{n=0}^G {G\choose n}(.1)^n(.9)^{G-n}
\sqrt{\frac nG\left(1-\frac nG\right)}>0.
\label{eq:kappa}
\end{equation}
Constant-score groups contribute zero in this limit. Equivalently, these numbers follow from zero advantages on such groups and population standardization on mixed groups; no division by zero is performed. The limiting derivatives are
\begin{equation}
\dot J_c=-.02\kappa_G<0,\qquad
\dot J_r=\frac7{150}\kappa_G>0.
\label{eq:gspodots}
\end{equation}
\begin{center}
\renewcommand{\arraystretch}{1.15}
\begin{tabular}{rrrr}
\toprule
$G$ & $\kappa_G$ & Standardized-limit $\dot J_c$ & Centered-only $\dot J_c$\\
\midrule
2 & $.0900000000$ & $-.0018000000$ & $-.0003000$\\
8 & $.2094920147$ & $-.0041898403$ & $-.0005250$\\
16 & $.2563031126$ & $-.0051260623$ & $-.0005625$\\
\bottomrule
\end{tabular}
\end{center}

Equal positive trajectory lengths add a common inverse-length factor to the update. Unequal lengths, nonindependent sampling, later clipped updates, or an additional loss gradient need not preserve the collinearity in Eq.~\eqref{eq:family_gspo_collinear}.

\subsection{Optimizer and parameterization boundaries}
\label{sec:family_coordinates}
For a positive diagonal preconditioner
$D=\diag(w_1,w_2,w_3)$, Eq.~\eqref{eq:family_dot} becomes
\begin{equation}
 g_c^\top Dg_r=\frac{ad^2}{m}
 \left[w_1am-w_2b^2+w_3d(1-d)\right].
 \label{eq:family_preconditioner}
\end{equation}
The sign is therefore not invariant under arbitrary parameter scaling.
At the original point, $D=\diag(2,1,1)$ changes the derivative from
$-.0006$ to $+.0003$.
More generally, a local parametrization $z=z(\theta)$ with Jacobian $J$
induces $JJ^\top g_r$ under an ordinary $\theta$-gradient step, giving success
derivative $g_c^\top JJ^\top g_r$, not generally $g_c^\top g_r$.
An orthonormal change of coordinates on the logit tangent space preserves
the inner product; an arbitrary rescaling or neural-network parametrization
does not.

At every interior point, $g_r$ has signs $(+,-,-)$. Idealized first-step
Adam applied to the full expected gradient with zero moments and negligible
optimizer epsilon consequently gives success derivative
$g_c^\top(1,-1,-1)=2ad>0$.
At the original point this derivative is $.06$. This is not the expectation of stochastic Adam, since its nonlinear
preconditioning does not commute with averaging.
A fixed-score exponential tilt
$p_\eta(i)\propto p_i\exp(\eta r_i)$ preserves the low-score conditional
proportions and hence exact calibration. Its initial success derivative is
\begin{equation}
 \operatorname{Cov}_p(c,r)=\operatorname{Var}_p(r)
 =a(1-a)(1-q)^2=\frac{ad^2}{1-a}>0.
 \label{eq:family_tilt}
\end{equation}
At the original point the tilt derivative is $.01$. These alternatives delimit an ordinary logit-gradient
claim; the family does not predict the actual optimizer trajectory of the
source training system.

\section{ProCredit: a score guarantee and its combined advantage}
\label{sec:procredit}

\subsection{Published score and credit rules}
ProCredit defines progress $\Phi_t\in[0,1]$ as the fraction of executable task checks passed after turn $t$ \citep[Section 3, PDF pp.~5--6]{ma2026procredit}. To avoid confusing correctness with the source's progress coefficient $c$, we write that coefficient as $b>0$. With outcome $R\in\{0,1\}$,
\begin{align}
r_t&=b(\Phi_t-\Phi_{t-1}),\nonumber\\
S&=R+\sum_{t=1}^T r_t=R+b(\Phi_T-\Phi_0).
\label{eq:proscore}
\end{align}
For a common initial progress, bounded final progress, and $R=1\Rightarrow\Phi_T=1$, a sure-success policy attains the maximum possible expected score. Under its availability, the score-optimal and success-optimal policy sets coincide \citep[Appendix G, Proposition 3, PDF p.~23]{ma2026procredit}. This is a global statement about expected score, not every intermediate parameter update.

The source standardizes trajectory scores within a task group and adds centered turn-level returns:
\begin{align}
A_i^{\rm traj}&=\frac{S_i-\bar S}{s_G(S)+\epsilon_s},\nonumber\\
G_{i,t}^{(\gamma)}&=\gamma^{T_i-t}R_i+
 \sum_{u=t}^{T_i}\gamma^{u-t}r_{i,u},\nonumber\\
A_{i,t}^{\rm turn}&=G_{i,t}^{(\gamma)}-
 \frac1{|\mathcal U|}\sum_{(j,u)\in\mathcal U}G_{j,u}^{(\gamma)},\nonumber\\
A_{i,t}&=A_i^{\rm traj}+A_{i,t}^{\rm turn}.
\label{eq:proadv}
\end{align}
Here $\mathcal U$ contains all turns of the sampled trajectories for the task. The source uses $\gamma=1$, so
$G_{i,t}=R_i+b(\Phi_{i,T_i}-\Phi_{i,t-1})$.
Its clipped token-ratio objective averages over generated action tokens and excludes tool responses (source Eq.~2); no KL or entropy term is used. Reported defaults are $G=8$ and $b=.5$ (Table I1, PDF p.~24).

\subsection{One turn, one token, and exact pair enumeration}
Consider ten initially unsatisfied task checks and one turn. Three one-token actions pass ten, one, or zero checks. Set
\begin{equation}
p=(.02,.75,.23),\quad c=(1,0,0),\quad
\Phi_T=(1,.1,0),\quad\Phi_0=0,\quad b=.5.
\end{equation}
Then $S=(1.5,.05,0)$. The full categorical simplex contains a sure-success vertex that uniquely maximizes score. Finite softmax logits parameterize its interior; we evaluate the local update at $p$.

With one turn per trajectory, $G_{i,1}=S_i$, and Eq.~\eqref{eq:proadv} reduces exactly to
\begin{equation}
A_i=\frac{S_i-\bar S}{s_G(S)+\epsilon_s}+(S_i-\bar S).
\label{eq:proone}
\end{equation}
The source's token-mean objective is now the group mean. At the behavior policy, with inactive clipping and fixed sampled advantages, the gradient is $G^{-1}\sum_iA_i(e_i-p)$.

For $G=2$, use population standard deviation, $\epsilon_s=0$, and zero advantage on ties. In an ordered pair $(i,j)$, define $d_{ij}=S_i-S_j$. Then
\begin{equation}
A_i=\sgn(d_{ij})+\frac{d_{ij}}2=-A_j,
\qquad
g^{(i,j)}=\frac{A_i}{2}(e_i-e_j).
\end{equation}
Summing the nine pairs with probabilities $p_ip_j$ gives
\begin{equation}
g_i=p_i\sum_jp_j\left[\sgn(S_i-S_j)+\frac{S_i-S_j}{2}\right].
\end{equation}
The resulting gradients are
\begin{align}
g_c&=(.0196,-.015,-.0046),\nonumber\\
g_S&=(.02865,-.013125,-.015525),\nonumber\\
g&=(.033925,.1509375,-.1848625).
\end{align}
Their directional products are
\begin{align}
g_c^\top g_S&=.00082983>0,\nonumber\\
g_c^\top g&=-.000748765<0,\nonumber\\
g_S^\top g&=.001860886875>0.
\label{eq:prodots}
\end{align}
The raw-score gradient would improve success locally, but the expected combined-advantage direction reduces it while increasing mean score. The standardized trajectory term contributes $-.00116368$ to the success derivative; the centered turn term contributes $+.000414915$. The turn term mitigates the unfavorable standardized contribution. Both terms have the same sign within a pair; the reversal arises from their relative weighting across random groups.

\paragraph{Relation to existing normalization analyses.}
\citet{bereket2025uncalibrated} analyze overconfident probability prediction under proper scoring rewards; \citet{wang2026darkroom} studies amplification of policy-dependent prediction rewards in long-horizon agents. Both identify failures alleviated by removing standard-deviation normalization. The construction here fixes verified progress and examines ProCredit's specific advantage while preserving its global success-optimal policy set. It contrasts with Section~\ref{sec:calibration}, where mean centering leaves the unfavorable raw-gradient direction unchanged.

\subsection{Default group size and removing standardization}
For a group of size $G$, enumerate all count triples $n=(n_1,n_2,n_3)$ with $\sum_kn_k=G$. Their probabilities are
\begin{equation}
P(n)=\frac{G!}{n_1!n_2!n_3!}\prod_kp_k^{n_k}.
\end{equation}
For each triple, compute the mean, standard deviation, and advantages in Eq.~\eqref{eq:proone}, then accumulate
\begin{equation}
\E[g]=\sum_nP(n)\sum_k\frac{n_k}{G}A_k(n)(e_k-p).
\label{eq:counts}
\end{equation}
There are $\binom{G+2}{2}$ triples: 45 at the source's default $G=8$. At $b=.5$ and $\epsilon_s=10^{-6}$, the success derivatives are
\begin{center}
\begin{tabular}{lr}
\toprule
Standard-deviation convention & $g_c^\top\E[g]$, $G=8$\\
\midrule
Population & $-.000929983748$\\
Sample & $-.000823033653$\\
\bottomrule
\end{tabular}
\end{center}
The displayed source formula contains $\epsilon_s$ but does not fix its value or an unambiguous training standard-deviation convention. These are explicit sensitivity calculations, not a reconstruction of every implementation detail.

Removing the divisor from the trajectory term leaves two copies of $S_i-\bar S$. Exact mean-baseline enumeration gives
\begin{equation}
\E[g_{\rm no\,std}]=2\frac{G-1}{G}g_S.
\label{eq:pronostd}
\end{equation}
The success derivative is positive for this construction: $.00082983$ at $G=2$ and $.0014522025$ at $G=8$. This intervention repairs the constructed ProCredit reversal, whereas Eq.~\eqref{eq:sweno} shows that it does not repair the calibrated SWE-RM family in Eq.~\eqref{eq:family_q_region}. For that two-score family, both standardized and centered credit are positive multiples of $g_r$ by Eqs.~\eqref{eq:family_gspo_collinear} and \eqref{eq:sweno}. Its mismatch is already present in the raw score; the ProCredit example instead acquires its unfavorable direction in the conversion to finite-group credit.

For fixed positive $\epsilon_s$, the small score difference between the two failed actions vanishes as $b\to0$, and the direction becomes favorable again. In the fixed-score large-$G$ limit, Eq.~\eqref{eq:proone} yields
$[1/(\sigma_S+\epsilon_s)+1]g_S$, also favorable here. The example therefore does not assert reversal at arbitrary group sizes or coefficients. It neither refutes ProCredit's global score-optimum proposition nor reproduces AppWorld or ToolSandbox training. One turn also removes temporal-discount effects from this calculation.

\section{Reproduction}
Both programs use the Python standard library and run independently.
Run \texttt{python reproduce\_checks.py} in this folder to write \texttt{results.json}. It verifies:
\begin{itemize}
\item ARM-FM's two path returns and reversal thresholds;
\item exact calibration, the conditional covariance, and the original ordinary-logit example;
\item the GSPO limiting binomial expression against full count enumeration at $G=2,8,16$;
\item ProCredit's nine ordered pairs against an independent count representation, the 45 states at $G=8$, and the no-standardization comparison.
\end{itemize}
Rational probabilities and nonstandardized expressions use Python's \texttt{Fraction}; standardized calculations use floating-point square roots. The output preserves the individual $G=8$ count states and their probabilities, so the reported expectation can be inspected without Monte Carlo uncertainty.

Run \texttt{python} followed by \path{reproduce_calibration_family.py};
the output is \path{reproduce_calibration_family.results.json}. This program checks the family identities using exact rational arithmetic at 1,200 base policies, the strict negative-region formula at 6,000 numerical points, and fixed-score ambient directional derivatives by 80 central finite differences. It independently enumerates multinomial count vectors in 128 GSPO cases with $G\in\{2,3,8,16\}$, population and sample standard deviations, and stabilizers $10^{-6}$ and $.05$, comparing the results with Eq.~\eqref{eq:family_gspo_direction}. It also checks the fixed-evaluator threshold, the slice root, a rectangle contained in the negative region, and the coordinate examples. These numerical points use the fixed seed \texttt{20261021}. The analytical derivation establishes the identities over their stated domains; finite checks corroborate it and do not estimate prevalence in trained systems. This program does not rerun the ProCredit construction.

\bibliographystyle{plainnat}
\begingroup
\raggedright
\bibliography{references}
\endgroup
\end{document}
